\section*{致谢}

本书的大部分材料，其初稿写于 2018 年春在约翰斯·霍普金斯大学开设的一门为期
一学期的专题课；Qingci An、Thomas Brazelton、tsil clingman、Daniel
Fuentes-Keuthan、Aurel Malapani-Scala、Mona Merling、David Myers、Apurv
Nakade、Martina Rovelli 与 Xiyuan Wang 在课上提出的敏锐问题，使这份初稿
获益良多。

进一步的修订是在 2018 年于 MIT 举行、以\trInfinityCategory{}的模型无关理论为主题的
Talbot 工作营期间完成的；工作营由 Eva Belmont、Calista Bernard、Inbar
Klang、Morgan Opie 与 Sean Pohorence 组织，Haynes Miller 任指导教师。
Kevin Arlin、Timothy Campion、Kyle Ferendo、Daniel Fuentes-Keuthan、Joseph
Helfer、Paul Lessard、Lyne Moser、Emma Phillips、Nima Rasekh、Martina
Rovelli、Maru Sarazola、Matthew Weatherley、Jonathan Weinberger、Laura
Wells 与 Liang Ze Wong 在那次工作营上所作的引人入胜的报告令我们深受启发；
与 Elena Dimitriadis Bermejo、Luciana Basualdo Bonatto、Olivia Borghi、
Tai-Danae Bradley、Tejas Devanur、Aras Ergus、Matthew Feller、Sina
Hazratpour、Peter James、Zhulin Li、David Myers、Maximilien Péroux、
Mitchell Riley、Luis Scoccola、Brandon Shapiro、Pelle Steffens、Raffael
Stenzel、Paula Verdugo 与 Marco Vergura 的无数次讨论同样令我们受益。

John Bourke 提出了引理 2.1.11，它统一了关于\trAdjunction{}的 2-\trCategory{}理论的若干结果的
证明。Denis-Charles Cisinski 提请我们注意 André Joyal 的一个观察，即今
练习 4.2.iii。Omar Antolín-Camarena 向我们谈到了定理 4.4.12 的条件（iii）
——事实证明，这是刻画稳定\trInfinityCategory{}的最便捷途径，借之可在命题 6.3.16 中证明
它们装配成一个\trInfinityCosmos{}。Gabriel Drummond-Cole 鼓励我们对第 5 章作了一番
亟需的结构重组。Anna Marie Bohmann 建议了练习 8.2.i。tsil clingman 向我们
指出 Tom Leinster 的一个观察，它启发了练习 8.3.ii。第 9 章的表述因 Kevin
Arlin 的观察而大为改善——他促成了该章的重新组织——也因 David Myers 而
改善，命题 9.1.8 所载结果即由他最先陈述。关于 $(\infty,1)$-\trCategory{}之\trGroupoid{}核
（groupoid core）的材料，得益于与 Alexander Campbell 和 Yuri Sulyma 的
讨论——后者在博士论文工作期间独立（且更早）地发展了这部分材料的大半。
11.3 节的结果，受 Simon Henry 在\trHomotopy{}类型论电子系列研讨会 [52] 上一次报告的
启发。Timothy Campion、Yuri Sulyma 与 Dimitri Zaganidis 各自就最终从本书
定稿中删去的材料提出了极好的建议。

Alexander Campbell 在附录 A 的写作期间数次受我们咨询，特别是关于换富集%
\trFunctor{}与底范畴之间微妙的相互作用。他还提供了命题 E.3.3 中 \emph{n}-\trQuasiCategory
之\trInfinityCosmos{}的证明。Naruki Masuda 抓出了附录 D 中的若干笔误与不一致，而
Christian Sattler 为命题 D.5.6 中的蕴涵（iii）$\Rightarrow$（i）给出了
一个更简单的证明。

tsil clingman 建议了一种在\trVirtualEquipment{}中排置前箭头（proarrow）
的更美观的方式，使第 8 章的展示图大为改观。Doug Ravenel 推广了以字符
「$\hookrightarrow$」表示 \trYonedaEmbedding{}（Yoneda embedding）的做法。Anna
Marie Bohmann、John Bourke、Alexander Campbell、Antoine Chambert-Loir、
Arun Debray、Gabriel Drummond-Cole、David Farrell、Harry Gindi、Philip
Hackney、Peter Haine、Dodam Ih、Stephen Lack、Chen-wei (Milton) Lin、
Naruki Masuda、Mark Myers、Viktoriya Ozornova、Jean Kyung Park、Emma
Phillips、Maru Sarazola、Yuri Sulyma、Paula Verdugo、Mira Wattal、Jonathan
Weinberger 与 Hu Xiao 曾来信指出笔误。

Peter May 建议了「\trInfinityCosmos」这一名称——较之先前的非正式用语，这是一大改进。
Mike Hopkins 提议了本书的书名。我们感谢剑桥大学出版社 Kaitlin Leach 与
Amy He 的悉心督导、Suresh Kumar 的技术支持、Anbumani Selvam 卓越的文字
编辑与 Dhivya Elavazhagan 的制作管理，并感谢匿名审稿人提出的敏锐建议。

最后，作者感谢国家科学基金会（资助号 DMS-1551129 与 DMS-1652600）、
澳大利亚研究理事会（资助号 DP160101519 与 DP190102432）以及约翰斯·霍普
金斯大学「催化剂」与「校长前沿奖」计划提供的经费支持。我们还要感谢约翰斯·
霍普金斯大学数学系、麦考瑞大学澳大利亚范畴论中心与数学科学研究所接待我们
各自的到访。而最要紧的，是感谢我们的家人给予的爱、支持、耐心与体谅——
没有这一切，本书绝无问世的可能。
